\section{Typed Plans and Verified Query Planning}
\label{sec-problem}

This section formalizes query plans as typed terms of the SPARQL
algebra and states the planning problem that \hermes solves.

\stitle{Plans}.
A plan is the logical form of a query, built one operator at a time.

\begin{definition}[Plan]
\label{def-plan}
A plan is a tuple $Q = (T, F, J, X, K, A,$ $H, R, \Pi, \delta, \ell, f)$:
a set $T$ of triple patterns;
a set $F$ of filter expressions;
a set $J$ of groups $(\omega, T')$ that combine a set $T'$ of triple
patterns with the plan by
$\omega \in \{\kw{optional}, \kw{union}, \kw{minus}\}$;
extensions $X$, pairs $(e, ?v)$ of an expression and a variable;
group keys $K \subseteq V$;
aggregates $A$, pairs $(\varphi, ?v)$ of an aggregate expression and a
variable;
group filters $H$;
order keys $R$;
projected variables $\Pi \subseteq V$;
a distinct flag $\delta$;
a slice $\ell$;
and a form $f \in \{\kw{SELECT}, \kw{ASK}\}$.
\end{definition}

A plan denotes an algebra term with a fixed operator order, written as
in relational algebra and applied right to left:
\begin{equation}
\label{eq-plan}
Q \;=\; \opslice_\ell\; \opdist\; \opproj_\Pi\; \opord_R\; \opsel_H\;
\opgrp_{K;A}\; \opsel_F\; \opext_X \bigl(\,{\bowtie}\,T \odot J\,\bigr),
\end{equation}
where ${\bowtie}\,T$ joins the triple patterns of $T$, $\odot J$ applies
each group of $J$ with its operator ($\ljoin$ for \kw{optional}, $\cup$
for \kw{union}, $-$ for \kw{minus}), an operator with an empty component
is the identity, and $\opdist$ applies only when the distinct flag is
set; an \kw{ASK} plan replaces the top by $\exists$.
This is the order of the SPARQL 1.1 algebra~\cite{sparql11}.
Rendering is deterministic, so a plan and its query are interchangeable,
and we write $\eval{Q}{\G}$ for its answers.
An \emph{action} $a$ is a total function from plans to plans that
applies one operator of Table~\ref{tab-ops}; the action \kw{remove}
retracts a triple pattern or a filter.
Plans are immutable, so a sequence of actions is a path
$Q_0 \to Q_1 \to \cdots \to Q_k$ from the empty plan $Q_0$.

\stitle{Typing}.
Let $\vars(\cdot)$ denote the variables of a pattern or expression.
A \emph{typing context} $\Gamma$ maps each variable to a set of types.
The typing context $\Gamma_Q$ of a plan $Q$ over an ontology $O$ is
the least context closed under the following rules, applied to every
triple pattern of $T$ and $J$:
\begin{gather*}
\frac{(?x,\ \kw{rdf{:}type},\ c) \quad c \in C}{c \in \Gamma(?x)}\ (\textsc{T-Type})
\\[1ex]
\frac{(?x,\ p,\ o) \quad d \in \dom(p)}{d \in \Gamma(?x)}\ (\textsc{T-Dom})
\\[1ex]
\frac{(s,\ p,\ ?y) \quad r \in \rng(p)}{r \in \Gamma(?y)}\ (\textsc{T-Rng})
\end{gather*}
Types are compared up to $\sub$.

\begin{definition}[Well-typed plan]
\label{def-typed}
A plan $Q$ is \emph{ill typed} over $O$ if
(1) a literal occurs in subject position;
(2) a literal is the object of a property $p$ whose range $\rng(p)$
contains a class;
(3) for some variable, $\Gamma_Q$ holds both a type in $\Lit$ and a
class; or
(4) a class or property in a namespace that $O$ describes is not in
$C$ or $E$.
Otherwise $Q$ is \emph{well typed}.
\end{definition}

Two classes in $\Gamma_Q(?x)$ that are incomparable under $\sub$ are
reported as a warning, not an error, because RDF makes the open-world
assumption and $O$ need not declare disjointness.

\begin{definition}[Well-formed plan]
\label{def-formed}
Let $\beta(Q)$ be the variables bound by $Q$: the variables of $T$,
of the \kw{optional} and \kw{union} groups of $J$, and of the extensions
$X$; the variables of \kw{minus} groups are not bound.
Let $\beta^+(Q) = \beta(Q) \cup \{?v \mid (\varphi, ?v) \in A\}$.
A plan is \emph{well formed} if
$\vars(e) \subseteq \beta(Q)$ for every filter and extension $e$;
no extension binds a variable of $T$ or $J$;
$K \subseteq \beta(Q)$;
the variables of $H$, $R$, and $\Pi$ are in $\beta^+(Q)$; and,
when $A$ or $K$ is non-empty, $\Pi \cap \beta(Q) \subseteq K$.
\end{definition}

A plan that is both well typed and well formed over $O$ is
\emph{valid}, written $O \entails Q$.
Validity is decided from $Q$ and $O$ alone, without evaluating $Q$
on $\G$.

\stitle{Goals}.
A question constrains the shape of its answer before any query exists:
``which department'' asks for entities of one class, ``how many'' for a
number, ``is there'' for a Boolean.

\begin{definition}[Goal]
\label{def-goal}
A goal is $g = (\vec{a}, \kappa, c, m)$: answer variables
$\vec{a}$, a kind
$\kappa \in \{\kw{entity}, \kw{literal}, \kw{number}, \kw{boolean}\}$,
an optional type $c$, and a cardinality $m \in \{\kw{one}, \kw{many}\}$.
A valid plan $Q$ \emph{satisfies} $g$, written
$Q \models g$, if:
for $\kappa = \kw{boolean}$, $f = \kw{ASK}$;
otherwise $f = \kw{SELECT}$, $\vec{a} \subseteq \beta^+(Q)$, the
projection is exactly $\vec{a}$, a \kw{number} is witnessed by an
aggregate or a numeric datatype in $\Gamma_Q$, an \kw{entity} has no
type in $\Lit$, and, if $c$ is given, some type in $\Gamma_Q$ is
subsumed by $c$.
The cardinality $m = \kw{one}$ is checked after execution:
$|\eval{Q}{\G}| = 1$.
\end{definition}

The goal is derived from the question by the planner, so it can be
wrong; satisfaction guarantees that the plan has the shape its goal
asks for, not that the goal reflects the question.

\stitle{The planning problem}.
\begin{definition}[Verified query planning]
\label{def-problem}
Given a question $q$, a graph $\G$, its ontology $O$, and a step budget
$B$, a policy $\rho$ states a goal $g$ and applies actions
$a_1, \dots, a_k$ with $k \le B$, producing plans
$Q_1, \dots, Q_k$.
After each action it receives an observation
$o_i = (\mathit{ok}_i, n_i, n_{i-1}, \mathit{rows}_i, \mathit{err}_i)$:
whether $Q_i$ is valid and runs, its number of solutions
$n_i = |\eval{Q_i}{\G}|$, the previous count, a sample of rows, and
an error message.
An action whose plan is not valid is rejected and not executed.
The output $Q_k$ must satisfy $O \entails Q_k$ and
$Q_k \models g$; among such outputs, the objective is the set
$F_1$ of $\eval{Q_k}{\G}$ against the intended answers of $q$.
\end{definition}

\begin{example}
\label{exa-plan}
Take the question ``Which US city hosted the most Winter Olympics?''
over a small graph of Olympic host cities.
The planner first states the goal $g = (?city, \kw{entity},
\kw{ex{:}City}, \kw{one})$.
The query it builds is
\begin{center}
\footnotesize
\begin{tabular}{@{}l@{}}
\texttt{SELECT ?city WHERE \{}\\
\texttt{\ \ ?city ex:hosted ?e . \ ?city a ex:City .}\\
\texttt{\ \ ?city ex:country ex:USA .}\\
\texttt{\ \ ?e ex:season "Winter" . \}}\\
\texttt{GROUP BY ?city ORDER BY DESC(COUNT(?e))}\\
\texttt{LIMIT 1}
\end{tabular}
\end{center}
and, as a term of Equation~\eqref{eq-plan},
\[
\opslice_{1}\; \opproj_{?city}\; \opord_{\downarrow n}\;
\opgrp_{?city;\, n \leftarrow \kw{COUNT}(?e)}\,
\bigl(t_1 \bowtie t_2 \bowtie t_3 \bowtie t_4\bigr),
\]
where
\begin{align*}
t_1 &= (?city,\ \kw{ex{:}hosted},\ ?e),\\
t_2 &= (?city,\ \kw{a},\ \kw{ex{:}City}),\\
t_3 &= (?city,\ \kw{ex{:}country},\ \kw{ex{:}USA}),\\
t_4 &= (?e,\ \kw{ex{:}season},\ \text{``Winter''}).
\end{align*}
Figure~\ref{fig-plan} shows the plan as an operator tree, built bottom
up by ten actions, with the number of solutions the explorer observed
after each one.
Joining $t_3$ and $t_4$ cuts 19 solutions to 8 and then 4; grouping
leaves 3 cities; the slice leaves 1, Lake Placid.
The verifier keeps two goals open until the end: $?city : \kw{ex{:}City}$
until $t_2$ states the type, since this graph declares no domain for
\kw{ex{:}hosted}, and a single projected answer until the slice and the
projection.
\kw{finish} is accepted only after the last action.
\end{example}

\begin{figure}[t]
\centering
\begin{tikzpicture}[
    op/.style={font=\small, inner sep=1.5pt},
    leaf/.style={font=\footnotesize, inner sep=1.5pt},
    count/.style={font=\scriptsize, text=hermesblue, anchor=west},
    act/.style={font=\scriptsize, text=mutedink, anchor=east},
    edge/.style={draw=mutedink, line width=0.5pt},
    node distance=0.42cm]
  \node[op] (slice) {$\opslice_{1}$};
  \node[op, below=of slice] (proj) {$\opproj_{?city}$};
  \node[op, below=of proj] (ord) {$\opord_{\downarrow n}$};
  \node[op, below=of ord] (grp) {$\opgrp_{?city;\, n \leftarrow \kw{COUNT}(?e)}$};
  \node[op, below=of grp] (j3) {$\bowtie$};
  \node[op, below left=0.42cm and 0.55cm of j3] (j2) {$\bowtie$};
  \node[op, below left=0.42cm and 0.55cm of j2] (j1) {$\bowtie$};
  \node[leaf, below left=0.42cm and 0.25cm of j1] (t1) {$t_1$\,\textsf{hosted}};
  \node[leaf, below right=0.42cm and 0.25cm of j1] (t2) {$t_2$\,\textsf{a City}};
  \node[leaf, below right=0.42cm and 0.25cm of j2] (t3) {$t_3$\,\textsf{country USA}};
  \node[leaf, below right=0.42cm and 0.25cm of j3] (t4) {$t_4$\,\textsf{season Winter}};
  \foreach \a/\b in {slice/proj, proj/ord, ord/grp, grp/j3, j3/j2, j3/t4, j2/j1, j2/t3, j1/t1, j1/t2}
    \draw[edge] (\a) -- (\b);
  \node[count] at ([xshift=3pt]slice.east) {1};
  \node[count] at ([xshift=3pt]proj.east) {1};
  \node[count] at ([xshift=3pt]ord.east) {3};
  \node[count] at ([xshift=3pt]grp.east) {3};
  \node[count] at ([xshift=3pt]j3.east) {4};
  \node[count] at ([xshift=3pt]j2.east) {8};
  \node[count] at ([xshift=3pt]j1.east) {19};
  \node[count] at ([xshift=3pt]t1.east) {19};
  \node[act] at ([xshift=-3pt]slice.west) {$a_9$};
  \node[act] at ([xshift=-3pt]proj.west) {$a_{10}$};
  \node[act] at ([xshift=-3pt]ord.west) {$a_8$};
  \node[act] at ([xshift=-3pt]grp.west) {$a_6, a_7$};
  \node[act] at ([xshift=-3pt]j3.west) {$a_5$};
  \node[act] at ([xshift=-3pt]j2.west) {$a_4$};
  \node[act] at ([xshift=-3pt]j1.west) {$a_3$};
  \node[act] at ([xshift=-3pt]t1.west) {$a_2$};
\end{tikzpicture}
\caption{The plan of Example~\ref{exa-plan} as an operator tree.
  Gray: the action that added each operator ($a_1$ states the goal).
  Blue: solutions observed after that action.}
\label{fig-plan}
\end{figure}

The intended answers are unknown during planning, so the objective
cannot be optimized directly; the observations are the evidence the
policy has.
Three properties make this action space suitable for planning.

\begin{proposition}[Monotone refinement]
\label{prop-monotone}
Let $Q$ be a plan with $J = A = \emptyset$ and no slice, and let
$Q'$ apply \kw{triple} or \kw{filter} to $Q$.
Then every mapping in $\eval{Q'}{\G}$, restricted to $\beta(Q)$, is
a mapping in $\eval{Q}{\G}$.
\end{proposition}

\begin{proof}[Proof sketch]
A mapping of the join of $Q$ with a triple pattern is the union of
compatible mappings of $Q$ and of the pattern, so its restriction to
$\beta(Q)$ is a mapping of $Q$.
Filtering only removes mappings.
\end{proof}

Proposition~\ref{prop-monotone} makes the count in an observation
informative: within the conjunctive fragment, a new constraint can only
narrow the answers seen through the variables already present, so a
drop from 53 solutions to 1 (Example~\ref{exa-motivation}) is evidence
that the constraint is selective.

\begin{proposition}[Static soundness]
\label{prop-sound}
If the verifier of Section~\ref{sec-framework} accepts $Q$, then
$O \entails Q$.
\end{proposition}

\begin{proposition}[Checking cost]
\label{prop-cost}
After a one-time closure of $\sub$, deciding $O \entails Q$ and
$Q \models g$ takes time linear in the size of $Q$, up to the
cost of set lookups in $O$.
\end{proposition}

Proofs of Propositions~\ref{prop-sound} and~\ref{prop-cost} are in
Appendix~A.
